\scanpage{7}
\uline{Lemme} Soit $Z$ un espace, $Z'$ et $Z''$ deux parties fermées, telles que
$Z'\cap Z''=\{c\}$.
Pour que $c$ ait un voisinage conique dans $Z$, il f. et s. qu'il ait vois. conique dans $Z'$ et $Z''$.

Il est donc suffisant, pour avoir \deleted{\unsure{une}} décomposition de la ficelle en deux, d'imposer la condition

Li c) \uline{$C$ est régulièrement immergé dans $U$.}

\uline{Exemple} Pour $t$ assez petit, les $\unsure{\dot U|\dot V}_{A,t}=S_{A,t}$ sont des \deleted{\unsure{les}} lieux de la ficelle.

\margin{\uline{NB} Pour que $C$ partie fermée de $X$,
$\subset U=X\setminus(A\cup B)$, soit un lieu,
il f. et s. que \deleted{\unsure{dans}} $C\ne\varnothing$ et que
$\unsure{(X/_{A,B,C},\{a,b\}),\{c\}}$ soit un lieu.}

\margin{\deleted{Si $D$ est une partie fermée}
Si $D$ est une partie fermée de $X$ telle que \deleted{$\unsure{D\subset X\setminus(A\cup B\cup C)}$}
$D\subset\unsure{U|X}_{A,C}$ \uline{ou} $D\subset U_{B,C}$
(p.ex. \unsure{$D$ connexe}) alors
$C$ est \unsure{encore} un lieu de la ficelle
$\unsure{(X/D,\{A,B\})}$.}

\uline{Lieux disjoints.}
Soit $F=(X,\{A,B\})$ une ficelle, $\mathcal L(F)$ l'ens. de ses lieux.
Je vais définir une relation symétrique

``$C$ et $D$ sont disjoints''. \qquad $(\Longrightarrow(C\ne D))$

Quand $C$ \uline{ou} $D$ sont $\in\{A,B\}$, alors
``disjoints'' signifie simplement \uline{distinct}
($\Longleftrightarrow$ disjoints au sens ens., $C\cap D=\varnothing$).

\margin{\unsure{ou bien} $\{C,D\}=\{A,B\}$,
\unsure{ou bien l'un des deux} $C,D$ $\in\{A,B\}$, l'autre}

\rotatebox{90}{\reflectbox{\unsure{EXTRA STRONG}}}
\rotatebox{90}{\reflectbox{\unsure{EXTRA STRONG}}}
\mspage{6}
